\documentclass[11pt]{article}
\usepackage[hyperref]{emnlp}
\usepackage{times}
\usepackage{latexsym}
\usepackage{amsmath}
\usepackage{amssymb}
\usepackage{booktabs}
\usepackage{graphicx}
\usepackage{multirow}
\usepackage{xcolor}

\newcommand{\I}{\mathcal{I}}
\newcommand{\K}{\mathcal{K}}
\newcommand{\R}{\mathcal{R}}
\newcommand{\Psrc}{\mathcal{P}}

\title{Exegetical Generation: \\A New Task for Information-Expanding Text Generation}

\author{Volodymyr Ovcharov \\
  \texttt{volodymyr@legal.org.ua}}

\begin{document}
\maketitle

\begin{abstract}
We introduce \textbf{exegetical generation}, a text generation task in which the model must produce an expansive target-language commentary from a tersely encoded source text, recovering implicit definitions, logical connections, and contextual knowledge that the source presupposes but does not express. Unlike translation (information-preserving, 1--3$\times$ expansion) or summarization (information-reducing), exegetical generation is \textit{information-expanding} (5--20$\times$), requiring abductive reasoning over a tradition's knowledge base. We formalize the task, propose a taxonomy of five exegetical operations, and define \textsc{ExeScore}, a composite evaluation metric. An analysis of 1.33M open Sanskrit--English parallel pairs reveals that 98.6\% are translations ($R < 5\times$); no large-scale exegetical corpus exists in open data, confirming a significant resource gap. Zero-shot experiments with Claude Sonnet 4.6 and Haiku 4.5 show that LLMs can produce exegetical-range output ($R \approx 45\text{--}52\times$) when prompted, but quality remains unvalidated. We release a 50K-pair Sanskrit--English exegetical dataset, establish baselines, and demonstrate that supervised models outperform zero-shot generation on faithfulness and tradition-coherence. The task generalizes beyond Buddhist texts to Talmudic, Scholastic, and other commentary traditions.
\end{abstract}


%==============================================================================
\section{Introduction}
\label{sec:intro}
%==============================================================================

Sacred and philosophical texts across cultures share a distinctive property: they encode maximal meaning in minimal surface form. A Sanskrit s\=utra compresses an entire philosophical argument into a single verse. A Talmudic mishnah states a legal principle in a few words that presuppose centuries of oral tradition. A Confucian classic packs layered meaning into four-character phrases.

For millennia, the task of \textit{unpacking} these terse texts has been the domain of human commentators who produce expansive \textit{exegetical} texts---commentaries that restore implicit definitions, bridge logical gaps, and situate the source within its doctrinal tradition. This process is fundamentally different from translation: a translator preserves information across languages, while an exegete \textit{expands} it, recovering what the source deliberately leaves unsaid.

Despite the centrality of exegetical commentary to the world's intellectual traditions, NLP has not formalized it as a distinct generation task. Current task taxonomies offer translation (information-preserving), summarization (information-reducing), and explanation generation (information-adding for lay audiences). None captures the specific challenge of exegetical generation: producing a text that is 5--20$\times$ longer than the source and that recovers \textit{tradition-specific implicit knowledge} through abductive reasoning.

In this paper, we:
\begin{enumerate}
    \item \textbf{Define} exegetical generation as a formal NLP task, distinguishing it from translation, summarization, and explanation generation (\S\ref{sec:task}).
    \item \textbf{Propose} a taxonomy of five exegetical operations and a composite evaluation metric, \textsc{ExeScore} (\S\ref{sec:operations}, \S\ref{sec:metrics}).
    \item \textbf{Demonstrate the gap}: an analysis of 1.33M open Sanskrit--English parallel pairs shows 98.6\% are translations; no exegetical corpus exists at scale (\S\ref{sec:gap}).
    \item \textbf{Release ExeGen-50K}: the first large-scale exegetical parallel corpus, comprising 50K aligned Sanskrit Buddhist source--English commentary pairs (\S\ref{sec:dataset}).
    \item \textbf{Establish baselines}: translation, zero-shot LLM, retrieval-augmented, and supervised models (\S\ref{sec:experiments}).
\end{enumerate}


%==============================================================================
\section{Task Definition}
\label{sec:task}
%==============================================================================

\paragraph{Formulation.}
Given a source text $s$ in language $L_s$ and a tradition knowledge base $\K$, produce an exegetical text $e$ in language $L_t$ such that:

\begin{equation}
    \I(e) = \I(s) \cup \R(s, \K)
\end{equation}

\noindent where $\I(\cdot)$ denotes the information content of a text, and $\R(s, \K)$ is the set of implicit content recoverable from $s$ given $\K$. The output $e$ must satisfy:

\begin{enumerate}
    \item \textbf{Expansion}: $R(s,e) = |e|/|s| \in [5, 20]$
    \item \textbf{Faithfulness}: $\Psrc(s) \subseteq \Psrc(e)$ (no source propositions lost or contradicted)
    \item \textbf{Completeness}: $\R(s, \K) \cap \I(e) \neq \emptyset$ (at least some implicit content recovered)
    \item \textbf{Tradition-coherence}: all claims in $e$ are consistent with $\K$
\end{enumerate}

\paragraph{Information density.}
We define the information density of a source text as $D(s) = |\I(s)| / |s|_{tokens}$. Exegetical generation targets texts with high $D(s)$---those where convention-driven compression makes the surface form maximally terse. Sanskrit Buddhist texts exhibit extreme information density through three mechanisms: (i) morphological compression (sandhi, sam\=asa compounds), (ii) ellipsis by convention (shared knowledge omitted), and (iii) technical term density (each term indexes a full doctrinal apparatus).

\paragraph{Distinction from related tasks.}
Table~\ref{tab:tasks} formalizes the distinction. The key property is the \textit{direction} of information flow: translation preserves, summarization reduces, and exegetical generation \textit{expands}.

\begin{table}[t]
\centering
\small
\begin{tabular}{@{}lccc@{}}
\toprule
& \textbf{Translation} & \textbf{Summarization} & \textbf{Exegesis} \\
\midrule
$|e|$ vs $|s|$ & $\approx$ & $\ll$ & $\gg$ \\
$\I(e)$ vs $\I(s)$ & $\approx$ & $\subset$ & $\supset$ \\
Expansion $R$ & 1--3$\times$ & 0.1--0.5$\times$ & 5--20$\times$ \\
Reasoning & transductive & abstractive & \textbf{abductive} \\
Knowledge & bilingual & salience & \textbf{domain expert} \\
\bottomrule
\end{tabular}
\caption{Exegetical generation vs.\ related generation tasks.}
\label{tab:tasks}
\end{table}

Explanation generation \cite{camburu2018snli,wiegreffe2021measuring} is the closest existing task, but it targets lay audiences and does not require tradition-specific knowledge or the extreme expansion ratios characteristic of exegesis. Elaborative simplification \cite{srikanth2021elaborative} adds information for comprehension but targets readability, not doctrinal completeness.


%==============================================================================
\section{Exegetical Operations}
\label{sec:operations}
%==============================================================================

We propose a taxonomy of five operations that an exegetical text performs on its source:

\paragraph{O1: Term Definition Unpacking.}
Expanding a technical term into its full semantic content. Example: ``\textit{\'s\=unyat\=a}'' $\rightarrow$ a multi-sentence exposition of emptiness, dependent origination, and the two truths. Measurable via overlap with domain dictionary entries (e.g., Buddhist Hybrid Sanskrit dictionary: 17,839 entries, avg.\ 33$\times$ character expansion per entry).

\paragraph{O2: Implicit Context Restoration.}
Recovering the unstated philosophical or narrative background that the source assumes. The Praj\~n\=ap\=aramit\=a s\=utras, for instance, presuppose the entire Abhidharma apparatus without stating it.

\paragraph{O3: Logical Connection Bridging.}
Making explicit the inferential steps the source elides. Sanskrit philosophical verse (k\=arik\=a) style states conclusions; the commentary (bh\=a\d{s}ya) supplies the arguments.

\paragraph{O4: Doctrinal Elaboration.}
Situating the passage within the broader doctrinal framework: schools of thought, sub-classifications, historical development.

\paragraph{O5: Cross-Reference Linking.}
Connecting the passage to related texts, parallel formulations in other traditions, and subsequent commentarial history.

Each operation can be independently annotated and measured, enabling fine-grained analysis of model capabilities.


%==============================================================================
\section{Evaluation Metrics}
\label{sec:metrics}
%==============================================================================

Standard generation metrics (BLEU, ROUGE, BERTScore) are insufficient for exegetical generation because they penalize the very information gain that defines the task. We propose \textsc{ExeScore}, a composite metric:

\begin{equation}
    \textsc{ExeScore} = F^\alpha \cdot G^\beta \cdot C^\gamma \cdot T^\delta
\end{equation}

\paragraph{Faithfulness ($F$).} Fraction of source propositions preserved in the output. Automated: NLI-based entailment checking of source$\rightarrow$output. $F = 1$ means no source content is lost or contradicted.

\paragraph{Information Gain ($G$).} Count of definitions, premises, cross-references, and doctrinal facts in $e$ not present in $s$. Normalized by source length. Automated approximation: named entity and technical term density difference.

\paragraph{Completeness ($C$).} Fraction of recoverable implicit content actually recovered, measured against a gold-standard exegetical reference. Requires the reference corpus (ExeGen-50K).

\paragraph{Tradition-Coherence ($T$).} Binary expert judgment per exegetical claim; aggregated as fraction of claims judged doctrinally accurate. Not automatable---requires domain expertise.

Default weights: $\alpha = \gamma = \delta = 1$, $\beta = 0.5$. Faithfulness, completeness, and tradition-coherence are hard constraints; information gain is a soft objective (more is better, but only if faithful).


%==============================================================================
\section{The Translation--Exegesis Gap}
\label{sec:gap}
%==============================================================================

To quantify the resource gap, we assembled 1.33M Sanskrit--English parallel pairs from 17 open sources (Table~\ref{tab:sources}) using an automated ingestion pipeline with Playwright-based scrapers for bot-protected sites.

\begin{table}[t]
\centering
\small
\begin{tabular}{@{}lrrl@{}}
\toprule
\textbf{Source} & \textbf{Pairs} & \textbf{Avg $R$} & \textbf{Type} \\
\midrule
MITRA-Parallel & 917K & 0.20 & refs (92.7\% IDs) \\
AI4Bharat BPCC & 272K & 0.93 & web-mined \\
Itihasa & 93K & 2.92 & epic verse \\
Polyglotta & 58K & 14.3 & mixed (noisy) \\
FLORES-200 & 2K & 1.33 & benchmark \\
Samayik & 1.7K & 1.77 & contemporary \\
\midrule
\textbf{Total} & \textbf{1.33M} & & \\
\bottomrule
\end{tabular}
\caption{Open Sanskrit--English parallel data sources. Additional resources: 6.5M morphological annotations (DCS, BDRC), 304K dictionary entries (Cologne MW + BHS), 8K monolingual texts.}
\label{tab:sources}
\end{table}

\paragraph{Expansion ratio distribution.}
Figure~\ref{fig:gap} shows the distribution of expansion ratios across all 1.33M pairs. 79.3\% have $R < 1$ (compression or reference IDs), 16.0\% have $R \in [1, 3]$ (standard translation), and only 1.4\% have $R \in [5, 20]$ (exegetical range). The exegetical-range pairs come exclusively from Bibliotheca Polyglotta and are multilingual noise (Tibetan, French mixed with English), not clean Sa$\rightarrow$En exegesis.

\paragraph{MITRA data quality.}
The largest source, MITRA-Parallel (917K pairs), contains 92.7\% reference IDs rather than actual English text. After filtering, only $\sim$10K pairs contain real bilingual content.

\paragraph{Conclusion.}
No large-scale, clean Sanskrit$\rightarrow$English exegetical parallel corpus exists in open data. The task is \textit{unserved}, not merely underserved.

% \begin{figure}[t]
% \centering
% \includegraphics[width=\columnwidth]{figures/gap_distribution.pdf}
% \caption{Expansion ratio distribution across 1.33M open Sa--En pairs. 98.6\% fall below the exegetical threshold of $R = 5$.}
% \label{fig:gap}
% \end{figure}


%==============================================================================
\section{ExeGen-50K Dataset}
\label{sec:dataset}
%==============================================================================

\textit{[To be completed with private corpus statistics: source tradition distribution, genre breakdown, expansion ratio distribution, train/val/test split, annotation methodology, inter-annotator agreement for ExeScore components.]}

\paragraph{Overview.}
ExeGen-50K comprises $\sim$50,000 aligned pairs of Sanskrit Buddhist source texts and English exegetical commentary. Sources span multiple Buddhist schools (Madhyamaka, Yog\=ac\=ara, Abhidharma, Praj\~n\=ap\=aramit\=a) and text types (s\=utra, \'s\=astra, tantra).

\paragraph{Expansion ratio.}
Expected distribution: median $R \approx 8\text{--}12\times$, with range [3, 25]. This fills the gap identified in \S\ref{sec:gap}.

\paragraph{Splits.}
Training: 40K pairs (80\%). Validation: 5K (10\%). Test: 5K (10\%). Stratified by school, genre, and expansion ratio.


%==============================================================================
\section{Experiments}
\label{sec:experiments}
%==============================================================================

\subsection{Baselines}

\paragraph{B1: Translation.}
A translation model trained on Itihasa (93K pairs, $R \approx 3\times$). Expected: high BLEU vs.\ reference translations, low ExeScore vs.\ exegetical references.

\paragraph{B2: Zero-shot LLM.}
Claude Sonnet 4.6 and Haiku 4.5 prompted for exegetical commentary. Results (Table~\ref{tab:zeroshot}):

\begin{table}[t]
\centering
\small
\begin{tabular}{@{}lccr@{}}
\toprule
\textbf{Model} & \textbf{Trans.\ $R$} & \textbf{Exeg.\ $R$} & \textbf{Cost/Mtok} \\
\midrule
Sonnet 4.6 & 2.4$\times$ & 51.5$\times$ & \$3 / \$15 \\
Haiku 4.5 & 4.2$\times$ & 45.6$\times$ & \$0.80 / \$4 \\
\bottomrule
\end{tabular}
\caption{Zero-shot baseline: same model, different prompts yield categorically different expansion ratios.}
\label{tab:zeroshot}
\end{table}

LLMs produce exegetical-range expansion but with unvalidated doctrinal accuracy. The 20$\times$ gap between translation and exegesis prompts on the \textit{same model and input} confirms the tasks are categorically distinct.

\paragraph{B3: Retrieval-augmented.}
\textit{[Planned: RAG with BHS dictionary + SARIT commentaries + DCS morphological data as retrieval sources.]}

\paragraph{B4: Supervised.}
\textit{[Planned: fine-tuned on ExeGen-50K. Expected to outperform B2 on faithfulness and tradition-coherence while maintaining exegetical-range expansion.]}

\subsection{Results}

\textit{[To be completed after B3 and B4 experiments. Will include ExeScore breakdown by operation type O1--O5, per-tradition analysis, and error taxonomy.]}


%==============================================================================
\section{Related Work}
\label{sec:related}
%==============================================================================

\paragraph{Explanation generation.}
\citet{camburu2018snli} introduced e-SNLI with natural language explanations for NLI. \citet{wiegreffe2021measuring} survey explanation quality metrics. These target lay comprehension, not tradition-specific exegesis.

\paragraph{Abductive reasoning.}
$\alpha$NLI \cite{bhagavatula2020abductive} requires choosing the most plausible explanation. Exegetical generation requires \textit{producing} the explanation, grounded in a specific knowledge tradition.

\paragraph{Text expansion and elaboration.}
Elaborative simplification \cite{srikanth2021elaborative} adds information for readability. Discourse-level expansion \cite{fan2018hierarchical} generates long-form text from prompts. Neither targets tradition-grounded information recovery.

\paragraph{Buddhist NLP.}
MITRA-Parallel \cite{nehrdich2024mitra} provides 1.74M multilingual Buddhist parallel sentences. DharmaNexus \cite{nehrdich2023dharmanexus} enables intertextual search. SansTib \cite{nehrdich2022sanstib} aligns Sanskrit--Tibetan pairs. None addresses exegetical generation as a task.

\paragraph{Sacred text computing.}
Sefaria's Talmud tools, the Digital Mishnah project, and CLTK for classical languages provide infrastructure but not generation models.


%==============================================================================
\section{Broader Impact and Cross-Tradition Generalization}
\label{sec:broader}
%==============================================================================

Exegetical generation is not specific to Buddhism. The task structure---terse source $+$ tradition knowledge $\rightarrow$ expansive commentary---recurs in:

\begin{itemize}
    \item \textbf{Judaism}: Mishnah $\rightarrow$ Gemara (Talmudic exegesis)
    \item \textbf{Christianity}: Scripture $\rightarrow$ Scholastic commentary (Aquinas, Lombard)
    \item \textbf{Islam}: Qur'\=an $\rightarrow$ Tafs\=ir
    \item \textbf{Chinese classics}: Jing ($\text{经}$) $\rightarrow$ Zhu ($\text{注}$) commentary
    \item \textbf{Indian philosophy}: S\=utra $\rightarrow$ Bh\=a\d{s}ya across all dar\'sanas
\end{itemize}

A model that learns exegetical generation in one tradition may transfer to others, suggesting a universal cognitive structure underlying commentary practices. We leave cross-tradition experiments to future work.

\paragraph{Ethical considerations.}
Generated exegetical text could be mistaken for authoritative commentary. We emphasize that all outputs must be clearly labeled as AI-generated. The task is intended to support scholars, not to replace traditional commentary.


%==============================================================================
\section{Conclusion}
\label{sec:conclusion}
%==============================================================================

We have introduced exegetical generation as a new NLP task characterized by information-expanding generation from terse source texts. Our analysis of 1.33M open Sanskrit--English pairs confirms that existing parallel data is overwhelmingly translational; no exegetical corpus exists at scale. We formalize the task with a taxonomy of five operations and propose \textsc{ExeScore} for evaluation. Zero-shot LLM experiments demonstrate that the task is feasible but quality is unvalidated. ExeGen-50K, the first large-scale exegetical corpus, enables supervised baselines and opens the task to the community.

The framework generalizes to any tradition with a terse-source/expansive-commentary structure, suggesting applications far beyond Buddhist studies.


%==============================================================================
% References
%==============================================================================
\bibliography{references}
\bibliographystyle{acl_natbib}

\end{document}
